%% kapa: exam=1 round=32 year=2021 period=1 form=A subject=economics subject_ko=경제학원론 seq=18 number=58
%% source: Q-Net 감정평가사 기출문제 (cst00309) artlSeq=5208397 file=extracted/1008/5208397/A형/2021년도 제32회 감정평가사 1차 1교시 A형.pdf
%% conversion: kordoc 4.24.51 -> tools/gen1.py (auto); review=proofread-sonnet; flags=none
\begin{question}{58}
우하향하는 장기평균비용에 관한 설명으로 옳은 것은? (단, 생산기술은 동일하다.)
\begin{choices}
  \item 생산량이 서로 다른 기업의 평균비용은 동일하다.
  \item 진입 장벽이 없는 경우 기업의 참여가 증가한다.
  \item 소규모 기업의 평균비용은 더 낮다.
  \item 장기적으로 시장에는 한 기업만이 존재하게 된다.
  \item 소규모 다품종을 생산하면 평균비용이 낮아진다.
\end{choices}
\end{question}
\answer{4}
